\documentclass[11pt]{article}

\usepackage[margin=1in]{geometry}
\usepackage{amsmath,amssymb,amsthm,mathtools}
\usepackage{booktabs}
\usepackage{enumitem}
\usepackage{hyperref}
\usepackage[T1]{fontenc}

\newtheorem{theorem}{Theorem}[section]
\newtheorem{proposition}[theorem]{Proposition}
\newtheorem{lemma}[theorem]{Lemma}
\newtheorem{remark}[theorem]{Remark}
\newcommand{\Z}{\mathbf Z}
\newcommand{\tensor}{\otimes}
\newcommand{\Spec}{\operatorname{Spec}}
\newcommand{\rank}{\operatorname{rank}}
\newcommand{\finrank}{\operatorname{finrank}}

\title{A Mathlib-Style Formalization of the Arbitrary-Prime\\
Grothendieck Power Counterexample:\\
Construction, Verified Core, and Exact Proof Boundary}
\author{Akhil Mathew and Codex}
\date{July 13, 2026}

\begin{document}
\maketitle

\begin{abstract}
This note records a new Lean 4 formalization of the arbitrary-prime construction.
For every prime $p$ it defines the intended Artinian base $R_p$ and the
coordinate algebra $A_p$, proves in Lean that $A_p$ is finite free of rank
$p^2$ over $R_p$, proves the defining equations, and proves that the carry
term which detects the $p^2$-power word is nonzero.  It also develops a
Mathlib-native convolution API which turns the expected skew-primitive
coproduct into the desired power-map calculation.

The file contains no \texttt{sorry}, \texttt{admit}, or axiom declaration.
It does not pretend that the whole arbitrary-prime Hopf calculation has
already been checked: the odd-prime divided-bridge closure identity and two
geometric-sum identities are represented by explicit certificate structures.
Once terms of those structures are constructed, the final theorem gives a
finite free commutative Hopf algebra of rank $p^2$ whose $p^2$-power map is
not the unit map.  Thus the formalization separates the fully verified
commutative-algebra and nonvanishing core from the remaining Hopf-polynomial
normalization problem.
\end{abstract}

\tableofcontents

\section{Files and reproducibility}

The main new Lean source is
\begin{center}
\texttt{FiniteFlatGroupSchemes/GrothendieckCounterexampleAllPrimes.lean}.
\end{center}
It is imported by the project root \texttt{FiniteFlatGroupSchemes.lean}.
The source uses Lean $4.31.0$ and Mathlib $v4.31.0$.

No second Mathlib checkout was downloaded into Dropbox.  The project entry
\texttt{.lake} is a symbolic link to
\begin{center}
\texttt{\$HOME/.cache/finite-flat-group-schemes/lake}.
\end{center}
The final verification command was
\begin{verbatim}
lake build
\end{verbatim}
and it completed successfully with $8561$ jobs in the already populated
external cache.  A direct bounded check of the new source also succeeded:
\begin{verbatim}
lake env lean \
  FiniteFlatGroupSchemes/GrothendieckCounterexampleAllPrimes.lean
\end{verbatim}

The implementation follows the style of the polished rank-four formalization
shared at
\begin{center}
\url{https://gist.github.com/j2d9w5xtjn-png/b1fc94332439e1d43944163e1dae3aef}.
\end{center}
In particular, it uses narrow imports, documented definitions, explicit
Mathlib algebra maps, and a theorem boundary visible in the types.  For
arbitrary $p$, however, a two-step \texttt{AdjoinRoot} presentation is more
natural than the four-dimensional custom normal form used when $p=2$.

\section{The uniform mathematical construction}

Fix a prime $p$.  For $1\leq i\leq p-1$ put
\[
  w_i=\frac{1}{p}\binom pi\in\Z.
\]
It is convenient to set $w_0=0$ and to write
\[
  W_c(T)=\sum_{i=0}^{p-1}w_i c^{i-1}T^i,
\]
where the $i=0$ summand is zero.  In Lean the natural number coefficient is
\texttt{dividedChoose p i := p.choose i / p}; the theorem
\texttt{prime\_mul\_dividedChoose} verifies
\[
  p w_i=\binom pi\qquad(0<i<p).
\]

The base ring is
\begin{equation}\label{eq:base}
 R_p=\Z[a,b]/
 \bigl(a^{p+1},\ b^{2p-1},\ a^p b^{p-1}+p\bigr).
\end{equation}
The intended coordinate algebra is
\begin{equation}\label{eq:A}
 A_p=R_p[U,V]/(F,G),
\end{equation}
where
\begin{align}
 G(V)&=V^p-a^pW_b(V),\label{eq:G}\\
 F(U,V)&=U^p-a b^{p-1}W_a(U)+b^pW_b(V).
 \label{eq:F}
\end{align}

Rather than formalize a simultaneous two-variable quotient and then run a
Gr\"obner-basis argument, the Lean file constructs
\[
  B_p=\operatorname{AdjoinRoot}_{R_p}(G),\qquad
  A_p=\operatorname{AdjoinRoot}_{B_p}(F).
\]
Both polynomials are monic of degree $p$.  This presentation makes finite
freeness a structural consequence of Mathlib's canonical power basis.

\section{The verified base-ring identities}

The variables in $\Z[a,b]$ are implemented as an
\texttt{MvPolynomial (Fin 2) \(\mathbb Z\)}.  The quotient classes are
\texttt{a p} and \texttt{b p}.  Lean checks the three defining relations
directly from membership in \texttt{Ideal.span}.  It also checks that the
base is nonzero by evaluating $a$ and $b$ at zero in $\Z/p$.

The following consequences are formal theorems:
\begin{align*}
 pa&=0,\\
 pa^p&=0,\\
 p^2&=0,\\
 pb^p&=0.
\end{align*}
They occur as \texttt{p\_mul\_a}, \texttt{p\_mul\_a\_pow},
\texttt{p\_squared\_eq\_zero}, and \texttt{p\_mul\_b\_pow}.

The subtle assertion is that one power earlier does not vanish:
\begin{equation}\label{eq:pbnonzero}
 p b^{p-1}\neq0\quad\text{in }R_p.
\end{equation}
The theorem \texttt{p\_mul\_b\_pred\_ne\_zero} proves
\eqref{eq:pbnonzero} uniformly in $p$.

\subsection{The coefficient detector}

For a polynomial $q\in\Z[a,b]$, define
\begin{equation}\label{eq:detector}
 \ell(q)=
 [b^{p-1}]q-p[a^pb^{2p-2}]q\in\Z/p^2.
\end{equation}
This is \texttt{topCoefficientDetector}.  The proof verifies separately that
$\ell$ annihilates an arbitrary polynomial multiple of each generator of the
ideal in \eqref{eq:base}.

For the monomial generators this follows from exponent bounds.  For the
inhomogeneous generator $h=a^pb^{p-1}+p$, one computes for arbitrary $r$:
\begin{align*}
 [b^{p-1}](rh)&=p[b^{p-1}]r,\\
 [a^pb^{2p-2}](rh)&=[b^{p-1}]r
   +p[a^pb^{2p-2}]r.
\end{align*}
Consequently
\[
 \ell(rh)=-p^2[a^pb^{2p-2}]r=0\quad\text{in }\Z/p^2.
\]
The finite three-generator span is unpacked with
\path{Ideal.mem_span_insert} and
\path{Ideal.mem_span_singleton'}; no polynomial normalization tactic is
trusted here.  Finally
\[
 \ell(pb^{p-1})=p\neq0\quad\text{in }\Z/p^2,
\]
where the last inequality is checked using \texttt{ZMod.val}.  Hence
$pb^{p-1}$ cannot lie in the defining ideal.

\section{Finite freeness and rank}

The helper \texttt{monicPolynomial} constructs a polynomial
\[
 X^p+\sum_{i<p}c_iX^i.
\]
Theorems \path{monicPolynomial_isMonicOfDegree},
\path{monicPolynomial_monic}, and
\path{monicPolynomial_natDegree} prove that it is monic of natural degree
$p$.  These lemmas are then applied to $G$ and $F$.

Mathlib supplies a power basis for the root of a monic polynomial.  The file
therefore constructs instances
\[
 \operatorname{Module.Free}_{R_p}(B_p),\quad
 \operatorname{Module.Finite}_{R_p}(B_p),\quad
 \operatorname{Module.Free}_{B_p}(A_p),\quad
 \operatorname{Module.Finite}_{B_p}(A_p).
\]
Transitivity supplies finite freeness of $A_p$ over $R_p$.  The checked rank
statements are
\begin{align*}
 \finrank_{R_p}B_p&=p,\\
 \finrank_{B_p}A_p&=p,\\
 \finrank_{R_p}A_p&=p^2.
\end{align*}
The last is the theorem \texttt{finrank\_A}.  Conceptually the basis is
\begin{equation}\label{eq:basis}
 \{U^iV^j:0\leq i,j<p\}.
\end{equation}
Lean also proves the two root relations \eqref{eq:G} and \eqref{eq:F} as
\texttt{V\_relation}, \texttt{v\_relation}, and
\texttt{U\_relation}.

\section{The nonzero carry in the coordinate algebra}

A small general lemma,
\texttt{scalar\_mul\_adjoinRoot\_pow\_ne\_zero}, says that if $f$ is monic,
$c\neq0$, and $i<\deg f$, then
\[
  c\,\overline X^i\neq0\quad\text{in }\operatorname{AdjoinRoot}(f).
\]
Importantly, the coefficient ring need not be a domain.  The proof applies a
coordinate of the canonical power-basis representation.

Apply this lemma first to $B_p/R_p$ with coefficient $pb^{p-1}$ and basis
vector $V^{p-1}$, then to $A_p/B_p$ with basis vector $U$.  Together with
\eqref{eq:pbnonzero}, this gives the fully checked theorem
\texttt{carry\_ne\_zero}:
\begin{equation}\label{eq:carry}
  p b^{p-1}U V^{p-1}\neq0\quad\text{in }A_p.
\end{equation}
This is the exact term which should occur as the image of $U$ under the
$p^2$-power word.

\section{Expected Hopf structure}

Put
\[
 L=1+aU,\qquad M=1+bV,\qquad \lambda=LM^{-1}.
\]
The nilpotence relations imply that $L$ and $M$ are units.  The expected
coproduct is
\begin{align}
 \Delta(U)&=U\tensor1+\lambda\tensor U,\label{eq:deltaU}\\
 \Delta(V)&=V\tensor\lambda^{-1}+1\tensor V,\label{eq:deltaV}\\
 \Delta(\lambda)&=\lambda\tensor\lambda.
 \label{eq:deltalambda}
\end{align}
The counit sends $U,V$ to zero and $\lambda$ to one.

For $p=2$, the existing file
\texttt{GrothendieckCounterexample.lean} proves closure by explicit
normalization in a four-dimensional algebra.  For odd $p$, substituting
\eqref{eq:deltaU}--\eqref{eq:deltaV} into \eqref{eq:F}--\eqref{eq:G}
requires the divided bridge, or truncated-log, identity from the accompanying
mathematical construction.  This is the genuinely new polynomial identity;
it is not a routine consequence of monicity.

\section{The exact Lean proof boundary}

The structure \texttt{HopfPackage p} records concrete data
\begin{itemize}[leftmargin=2em]
\item $\lambda$ and $\lambda^{-1}$, together with their inverse and
      $LM^{-1}$ identities;
\item an $R_p$-algebra homomorphism
      $\Delta:A_p\to A_p\tensor_{R_p}A_p$;
\item a counit and antipode;
\item coassociativity, the two counit laws, the two antipode laws;
\item formulas \eqref{eq:deltaU}--\eqref{eq:deltalambda} and their inverse
      analogues.
\end{itemize}
From a term $H:\texttt{HopfPackage p}$, the definitions
\texttt{H.toBialgebra} and \texttt{H.toHopfAlgebra} construct actual Mathlib
instances using \texttt{Bialgebra.ofAlgHom} and
\texttt{HopfAlgebra.ofAlgHom}.  Thus \texttt{HopfPackage} is not an axiom or
an opaque proposition asserting existence: a future proof must supply every
map and every law.

The only further arithmetic input is
\texttt{PowerCertificate p H}, containing the two identities
\begin{align}
 \sum_{i=0}^{p^2-1}\lambda^i
   &=p b^{p-1}V^{p-1},\label{eq:geom1}\\
 \sum_{i=0}^{p^2-1}\lambda^{-i}
   &=p b^{p-1}V^{p-1}.
 \label{eq:geom2}
\end{align}
Nonvanishing is no longer a field of this certificate; it is the independent
theorem \eqref{eq:carry}.

The present formalization therefore has the status summarized below.
\begin{center}
\begin{tabular}{@{}p{0.56\textwidth}p{0.29\textwidth}@{}}
\toprule
Claim & Lean status\\
\midrule
Definition and nontriviality of $R_p$ & checked\\
Relations $a^{p+1}=0$, $b^{2p-1}=0$, $a^pb^{p-1}+p=0$ & checked\\
$p^2=0$, $pb^p=0$, and $pb^{p-1}\neq0$ & checked\\
Definition of $B_p$ and $A_p$ by monic equations & checked\\
Finite free $A_p/R_p$ of rank $p^2$ & checked\\
$pb^{p-1}UV^{p-1}\neq0$ & checked\\
Artinianness of $R_p$ as a typeclass instance & not yet encoded\\
Closure of the proposed coproduct for arbitrary $p$ & explicit
  \texttt{HopfPackage} obligation\\
Geometric identities \eqref{eq:geom1}--\eqref{eq:geom2} & explicit
  \texttt{PowerCertificate} obligations\\
Final counterexample theorem & checked conditional on the two certificates\\
\bottomrule
\end{tabular}
\end{center}

\section{The convolution power calculation}

Assume $H:\texttt{HopfPackage p}$.  The file defines the universal point as
the identity algebra homomorphism in Mathlib's convolution monoid and defines
\texttt{powerMap p H n} as its $n$-th convolution power.  If
\[
 S_n(x)=1+x+\cdots+x^{n-1},
\]
then induction from the skew-primitive coproduct gives checked formulas
\begin{align}
 [n]^\#(U)&=S_n(\lambda)U,\label{eq:powerU}\\
 [n]^\#(V)&=V S_n(\lambda^{-1}),\label{eq:powerV}\\
 [n]^\#(\lambda)&=\lambda^n.
\end{align}
These are \texttt{powerMap\_U}, \texttt{powerMap\_v}, and
\texttt{powerMap\_lambda}.

Given \texttt{PowerCertificate p H}, equation \eqref{eq:powerU} yields
\[
 [p^2]^\#(U)=p b^{p-1}U V^{p-1}.
\]
By \eqref{eq:carry}, this is nonzero.  Hence the $p^2$-power word is not the
unit word.  Lean packages this as
\texttt{powerMap\_p\_squared\_U\_ne\_zero} and finally as
\texttt{counterexample}:
\[
 \finrank_{R_p}A_p=p^2,
 \qquad [p^2]^\#\neq e^\#.
\]

\section{Why the certificate boundary is useful}

There are two tempting but unsound shortcuts which the source avoids.
First, one could simply declare a Hopf algebra instance on $A_p$.  This would
hide the main odd-prime theorem in typeclass search.  Second, one could use an
axiom stating that the carry is nonzero.  The coefficient detector and power
bases now eliminate the second shortcut completely, while
\texttt{HopfPackage} exposes every field required to eliminate the first.

The boundary also permits the remaining calculation to be developed without
destabilizing the compiled commutative-algebra core.  A candidate coproduct
can be built with \texttt{AdjoinRoot.lift}; Lean will reduce its existence to
the two relation-preservation identities.  After those are proved,
coassociativity and the counit laws should be discharged by extensionality on
the two roots $U,V$.  The geometric identities can then be treated as a
separate truncated-binomial computation.

\section{Recommended next formalization steps}

The shortest route to a fully unconditional theorem appears to be:
\begin{enumerate}[leftmargin=2em]
\item Define finite geometric inverses for $L=1+aU$ and $M=1+bV$, proving in
      Lean that they are units from $a^{p+1}=0$ and $b^{2p-1}=0$.  This makes
      $\lambda=LM^{-1}$ unconditional.
\item State the divided-bridge identity as a standalone equality in a
      bounded polynomial quotient.  Prove it first before constructing any
      typeclass instances.
\item Use \texttt{AdjoinRoot.lift} twice to construct the coproduct from the
      two relation-preservation theorems.
\item Prove group-likeness of $\lambda$, coassociativity, counit, and antipode
      identities by \texttt{AdjoinRoot} extensionality.
\item Prove \eqref{eq:geom1} and derive \eqref{eq:geom2} by the inversion
      symmetry if possible.
\item Construct explicit terms of \texttt{HopfPackage p} and
      \texttt{PowerCertificate p H}; the already compiled theorem
      \texttt{counterexample} then closes the argument.
\item Finally provide an \texttt{IsArtinianRing (R p)} instance, most likely
      by exhibiting $R_p$ as a finite algebra over $\Z/p^2$.
\end{enumerate}

All computational experiments should remain bounded.  The present proof of
the top coefficient is symbolic and uses no Gr\"obner basis; compilation took
seconds and required no large-memory search.

\section{Conclusion}

The arbitrary-prime construction now has a substantial Mathlib-style formal
core.  For every prime $p$, Lean verifies the exact coordinate algebra, its
iterated power basis, finite freeness and rank $p^2$, the defining equations,
and the nonzero carry which witnesses failure of the $p^2$-power word.  The
convolution calculation from a skew-primitive Hopf structure to the final
counterexample is also verified.

What remains is sharply localized: the odd-prime divided-bridge closure
identity, the associated Hopf maps and laws, the two geometric sums, and an
Artinian typeclass instance for the base.  These are explicit theorem
obligations, not assumptions hidden behind \texttt{sorry} or an axiom.

\end{document}
